\chapter{Bewegte Frames und ihre Grenzen}

Das vorangehende Kapitel hat ein lokales Eisenbahn-Frame konstruiert, weil das
laufende Alignment eine Querschnittsorientierung erforderte. Dieses Kapitel
benennt die Abhängigkeiten jeder Richtung und die Grenze der
Erklärungskonvention.

\section{Longitudinalrichtung}

Für die realisierte räumliche Trajektorie \(\Gamma(s)\) lautet die
Einheitstangente
\[
\mathbf t(s)=\frac{\Gamma'(s)}{\|\Gamma'(s)\|}.
\]
Sie folgt aus horizontaler Pose, vertikalem Profil, metrischer Interpretation
und der Orientierung wachsenden \(s\). Horizontale Krümmung allein bestimmt
ihre dreidimensionale Richtung nicht.

\section{Nicht überhöhter Querschnitt}

Die Realisierung liefert einen verwendbaren Vertikalbezug
\(\mathbf z_{\mathrm{ref}}(s)\): Lokal kann dies die Vertikale einer ebenen
Engineering-Realisierung sein; bei einer Erdrealisierung können
Oberflächennormale und Gravitationskonvention beteiligt sein. Orthogonale
Projektion dieses Bezugs auf \(\mathbf t\), Normierung und rechtshändige
Vervollständigung ergeben eine mögliche Konstruktion von
\(F_0=(\mathbf t,\mathbf l_0,\mathbf z_0)\). Wo die Projektion schlecht
konditioniert ist, wird eine stetige Transportregel benötigt.

Diese Konstruktion bleibt bewusst qualifiziert. Die Thesis wählt kein
universelles Eisenbahn-Frame für jede Realisierung. Explizit bleiben müssen:
\begin{itemize}
\item Die Longitudinalrichtung folgt der realisierten Trajektorie und der
  Parameterorientierung.
\item Das nicht überhöhte Quer--Vertikal-Paar folgt einer angegebenen
  Vertikal- und Transportkonvention.
\item \(u(s)\) oder \(\phi(s)\) dreht dieses Paar um \(\mathbf t\).
\item Koordinatenkomponenten folgen der gewählten Darstellung.
\end{itemize}

\section{Warum ein reines Frenet-Frame nicht genügt}

Eine nur aus der Krümmung abgeleitete Frenet-Normale ist für
\(\kappa\rightarrow0\) nicht definiert oder numerisch instabil. Das laufende
Alignment enthält eine anfängliche Gerade; ein Eisenbahn-Frame muss seine
Orientierung daher stetig durch diesen Abschnitt tragen. Dies ist kein neues
architektonisches Objekt, sondern eine Anwendbarkeitsbedingung der
Frame-Auswertung.

\section{Bestimmtheitsgrade}

An einem festen \(s\) bleiben Alignment-Identität und geordnetes Gesetz
erhalten. Anfangspose und metrische Realisierung legen die ebene Platzierung
fest. Das Profil legt Höhe und Neigung fest, soweit definiert. Das
Überhöhungsgesetz legt den Roll nur unter einer Spurweiten- und
Vorzeichenkonvention fest. Eine verbleibende Gier- oder
Querschnittskonvention kann frei bleiben, bis die gewählte Frame-Regel sie
auflöst. Die Orientierung des Fahrzeugkastens bleibt frei, bis ein
Fahrzeugmodell und Betriebszustand bereitstehen.

\section{World--Track- und Vehicle-Space-Verwendung}

Die World--Track-Projektion konsumiert das qualifizierte Frame:
\(\mathbf t\) definiert die Longitudinalrichtung, \(\mathbf l_\phi\) den
Querabstand und \(\mathbf z_\phi\) die Querschnittshöhe. Eine
Vehicle-Space-Interpretation kann anschließend dieses Gleis-Frame konsumieren,
benötigt aber zusätzlich fahrzeugspezifische Daten. Das Frame bewahrt den
Alignment-Bezug und die Provenienz seiner Realisierung; es kann weder
entscheiden, ob ein Messpunkt zum beabsichtigten Objekt gehört, noch ob ein
Lichtraumergebnis akzeptabel ist.

\begin{summary}
Ein bewegtes Eisenbahn-Frame ist eine ausgewertete Beziehung zwischen
intrinsischer Geometrie, vertikaler Realisierung, Überhöhung und Konvention.
Seine Komponenten sind nur zusammen mit diesen Abhängigkeiten aussagekräftig.
\end{summary}
